% =========================================================
% Chapter: Sound Waves, Doppler Effect, Beats, Resonance, and Sonic Boom
% POSN 1 Physics Preparation Notes
% Language: Thai
% Compiler: XeLaTeX
%
% Required environments/packages from the main project:
%   amsmath, amssymb, siunitx, tikz, tcolorbox, graphicx
%   examplebox, exercisebox
% =========================================================

\chapter{เสียงและคลื่นเสียงในอากาศ: ดอปเพลอร์และโซนิกบูม}

\begin{center}
    {\Large \textbf{Sound Waves, Doppler Effect, and Sonic Boom}}\\[4pt]
    {\normalsize เอกสารเตรียมสอบคัดเลือก สอวน. ฟิสิกส์ ค่าย 1}\\[6pt]
    {\normalsize จัดทำโดย \textbf{Tames Thanakrit Damduan}}\\[2pt]
    {\footnotesize \textcopyright\ 2026 Tames Thanakrit Damduan. All rights reserved.}\\[8pt]
\end{center}

\sloppy

\section*{เป้าหมายของบทเรียน}

หลังจากเรียนบทนี้ นักเรียนควรทำสิ่งต่อไปนี้ได้

\begin{enumerate}
    \item อธิบายว่าเสียงในอากาศเป็นคลื่นกลตามยาวที่ประกอบด้วยบริเวณอัดและบริเวณขยายได้
    \item แยกความแตกต่างระหว่างอัตราเร็วของคลื่นเสียงกับความเร็วการสั่นของอนุภาคอากาศได้
    \item ใช้ความสัมพันธ์ $T=1/f$ และ $c_s=f\lambda$ แก้โจทย์พื้นฐานได้
    \item ใช้ความสัมพันธ์ของความเข้มเสียงและกฎกำลังสองผกผันได้
    \item ใช้สมการดอปเพลอร์เมื่อแหล่งกำเนิดหรือผู้สังเกตเคลื่อนที่ได้
    \item วิเคราะห์ความถี่สะท้อนจากวัตถุเคลื่อนที่และความถี่บีตได้
    \item ใช้เลขมัค มุมมัค และเรขาคณิตของโซนิกบูมได้
    \item แก้ชุดข้อสอบฝึกระดับคัดเลือกเข้า สอวน. ที่เชื่อมหลายแนวคิดได้
\end{enumerate}

\begin{tcolorbox}[title={ขอบเขตของบทนี้}]
บทนี้เน้นภาพทางกายภาพที่สังเกตได้และการแก้โจทย์ด้วยพีชคณิต ได้แก่ คลื่นตามยาว ความถี่ ความยาวคลื่น ความเข้มเสียง ดอปเพลอร์ บีต การสั่นพ้อง และโซนิกบูม

โจทย์ท้ายบทออกแบบให้ยากจากการเชื่อมหลายขั้นและการเลือกเครื่องหมายที่ถูกต้อง โดยไม่ต้องใช้แคลคูลัส
\end{tcolorbox}

% =========================================================
\section{เสียงคืออะไร}

\subsection{เสียงเป็นคลื่นกล}

คลื่นเสียงต้องอาศัยตัวกลางที่มีมวลและมีแรงคืนสภาพ อากาศทำหน้าที่ทั้งสองอย่าง

\begin{enumerate}
    \item โมเลกุลอากาศมีมวล จึงมีความเฉื่อย
    \item เมื่ออากาศถูกอัด ความดันเพิ่มขึ้นและพยายามดันให้กลับสู่สภาพเดิม
\end{enumerate}

การรบกวนจึงเดินทางต่อไปได้จากชั้นอากาศหนึ่งไปยังชั้นอากาศถัดไป

\begin{tcolorbox}[title={ข้อสรุปสำคัญ}]
เสียงไม่สามารถเดินทางในสุญญากาศ เพราะสุญญากาศไม่มีอนุภาคที่รับและส่งต่อการรบกวนเชิงกล
\end{tcolorbox}

\subsection{คลื่นตามยาวในอากาศ}

ให้ทิศการเดินทางของเสียงเป็นแนวแกน $x$ อนุภาคอากาศจะสั่นกลับไปกลับมาตามแนวเดียวกับแกน $x$ จึงเรียกว่า \textbf{คลื่นตามยาว}

บริเวณที่โมเลกุลอยู่ชิดกันกว่าปกติเรียกว่า \textbf{บริเวณอัด} ความดันและความหนาแน่นสูงกว่าค่าสมดุล

บริเวณที่โมเลกุลอยู่ห่างกันกว่าปกติเรียกว่า \textbf{บริเวณขยาย} ความดันและความหนาแน่นต่ำกว่าค่าสมดุล

\begin{center}
\begin{tikzpicture}[x=0.72cm,y=0.72cm]
    \draw[->,thick] (-0.5,0) -- (14.5,0) node[right] {$x$};
    \foreach \x in {0,0.28,0.56,1.8,2.5,3.2,4.4,4.68,4.96,6.2,6.9,7.6,8.8,9.08,9.36,10.6,11.3,12.0,13.2,13.48,13.76}{
        \fill (\x,0) circle (1.7pt);
    }
    \draw[decorate,decoration={brace,amplitude=5pt}] (0,-0.45) -- (0.56,-0.45)
      node[midway,below=6pt]{บริเวณอัด};
    \draw[decorate,decoration={brace,amplitude=5pt}] (1.8,-0.45) -- (3.2,-0.45)
      node[midway,below=6pt]{บริเวณขยาย};
    \draw[->,very thick] (5.5,1.0) -- (8.0,1.0) node[midway,above]{ทิศการเดินทางของคลื่น};
    \draw[<->,thick] (10.9,0.7) -- (11.7,0.7) node[midway,above]{การสั่นของอนุภาค};
\end{tikzpicture}
\end{center}

\subsection{อัตราเร็วคลื่นไม่ใช่ความเร็วอนุภาค}

ข้อผิดพลาดที่พบบ่อยคือคิดว่าโมเลกุลอากาศเคลื่อนที่จากลำโพงมาถึงหูด้วยอัตราเร็วประมาณ $340\,\si{m/s}$ ความจริงมีปริมาณสองชนิดที่ต้องแยกจากกัน

\begin{enumerate}
    \item $c_s$ คืออัตราเร็วที่รูปแบบการรบกวนหรือเฟสของคลื่นเดินทาง
    \item $u_p$ คือความเร็วขณะหนึ่งของอนุภาคอากาศซึ่งสั่นรอบตำแหน่งสมดุล
\end{enumerate}

โดยทั่วไปสำหรับเสียงความเข้มไม่สูงมาก

\[
|u_p|\ll c_s
\]

โมเลกุลหนึ่งเคลื่อนที่ไปกลับเป็นระยะสั้น ๆ แต่พลังงานและข้อมูลของคลื่นเดินทางได้ไกล

% =========================================================
\section{ปริมาณพื้นฐานของคลื่นเสียง}

\subsection{คาบ ความถี่ และความยาวคลื่น}

ถ้าแหล่งกำเนิดสั่นครบหนึ่งรอบในเวลา $T$ ปริมาณ $T$ เรียกว่า \textbf{คาบ} และจำนวนรอบต่อหนึ่งวินาทีเรียกว่า \textbf{ความถี่}

\[
\boxed{f=\frac{1}{T}}
\]

เมื่อคลื่นเดินทางด้วยอัตราเร็ว $c_s$ ในเวลาหนึ่งคาบ คลื่นจะเคลื่อนที่ไปได้หนึ่งความยาวคลื่น $\lambda$ ดังนั้น

\[
\lambda=c_sT
\]

แทน $T=1/f$

\begin{tcolorbox}[title={ความสัมพันธ์พื้นฐานของคลื่นเสียง}]
\[
\boxed{c_s=f\lambda}
\]
\end{tcolorbox}

สูตรนี้เชื่อมสามปริมาณ แต่ไม่ได้บอกว่าอนุภาคอากาศเคลื่อนที่ไปพร้อมกับคลื่น อนุภาคเพียงสั่นรอบตำแหน่งสมดุล ส่วนรูปแบบการอัดและการขยายเดินทางด้วยอัตราเร็ว $c_s$

\begin{examplebox}{ตัวอย่าง: ความยาวคลื่นและเวลาเดินทาง}
เสียงมีความถี่ $680\,\si{Hz}$ เดินทางในอากาศด้วยอัตราเร็ว $340\,\si{m/s}$

ความยาวคลื่นเป็น

\[
\lambda=\frac{c_s}{f}
=\frac{340}{680}
\]

\[
\boxed{\lambda=0.500\,\si{m}}
\]

ถ้าเสียงเดินทางเป็นระยะ $170\,\si{m}$ เวลาที่ใช้คือ

\[
t=\frac{s}{c_s}
=\frac{170}{340}
\]

\[
\boxed{t=0.500\,\si{s}}
\]
\end{examplebox}

\subsection{ความเข้มเสียงเบื้องต้น}

ความเข้มเสียง $I$ คือกำลังเสียง $\mathcal P$ ที่ผ่านพื้นที่ตั้งฉากหนึ่งหน่วย

\[
\boxed{I=\frac{\mathcal P}{A}}
\]

ถ้าแหล่งกำเนิดจุดแผ่เสียงเท่ากันทุกทิศ พลังงานกระจายบนผิวทรงกลมรัศมี $r$ ซึ่งมีพื้นที่ $4\pi r^2$ จึงได้

\begin{tcolorbox}[title={ความเข้มจากแหล่งกำเนิดจุด}]
\[
\boxed{I=\frac{\mathcal P}{4\pi r^2}}
\]

สำหรับจุดสองตำแหน่งจากแหล่งเดียวกัน

\[
\boxed{\frac{I_1}{I_2}=\left(\frac{r_2}{r_1}\right)^2}
\]
\end{tcolorbox}

ดังนั้นเมื่อระยะเพิ่มเป็นสองเท่า ความเข้มเหลือหนึ่งในสี่ ไม่ใช่หนึ่งในสอง

\begin{examplebox}{ตัวอย่าง: กฎกำลังสองผกผัน}
ผู้สังเกตคนหนึ่งอยู่ห่างลำโพง $5.0\,\si{m}$ แล้วเดินออกไปจนอยู่ห่าง $15.0\,\si{m}$ โดยถือว่าลำโพงเป็นแหล่งกำเนิดจุด จงหาอัตราส่วนความเข้มใหม่ต่อความเข้มเดิม

\[
\frac{I_{15}}{I_5}
=\left(\frac{5.0}{15.0}\right)^2
\]

\[
\boxed{\frac{I_{15}}{I_5}=\frac{1}{9}}
\]
\end{examplebox}

% =========================================================
\section{ปรากฏการณ์ดอปเพลอร์ของเสียง}

\subsection{แนวคิดจากระยะห่างระหว่างหน้าคลื่น}

ปรากฏการณ์ดอปเพลอร์คือการที่ความถี่ที่ผู้สังเกตได้รับแตกต่างจากความถี่ที่แหล่งกำเนิดปล่อยออกมา เพราะแหล่งกำเนิดหรือผู้สังเกตเคลื่อนที่เทียบกับตัวกลาง

ให้เสียงมีอัตราเร็ว $c_s$ เทียบกับอากาศ แหล่งกำเนิดปล่อยเสียงด้วยความถี่ $f_s$ และคาบ

\[
T_s=\frac{1}{f_s}
\]

ถ้าแหล่งกำเนิดอยู่นิ่ง ระหว่างการปล่อยยอดคลื่นสองยอด คลื่นยอดแรกเดินทางได้ระยะ

\[
\lambda=c_sT_s=\frac{c_s}{f_s}
\]

ถ้าแหล่งกำเนิดเคลื่อนที่เข้าหาผู้สังเกตด้วยอัตราเร็ว $v_s$ ในหนึ่งคาบแหล่งกำเนิดเคลื่อนที่ตามยอดคลื่นแรกไปเป็นระยะ $v_sT_s$ ระยะห่างระหว่างยอดคลื่นด้านหน้าจึงเป็น

\[
\lambda_{\rm front}
=c_sT_s-v_sT_s
=\frac{c_s-v_s}{f_s}
\]

เมื่อผู้สังเกตเคลื่อนที่เข้าหาแหล่งกำเนิดด้วยอัตราเร็ว $v_o$ ผู้สังเกตพบยอดคลื่นด้วยอัตราเร็วสัมพัทธ์ $c_s+v_o$ ดังนั้น

\[
f_o
=\frac{c_s+v_o}{\lambda_{\rm front}}
\]

แทนค่า $\lambda_{\rm front}$

\begin{tcolorbox}[title={สมการดอปเพลอร์ของเสียงในหนึ่งมิติ}]
\[
\boxed{
f_o=f_s\frac{c_s+v_o}{c_s-v_s}
}
\]

กำหนดเครื่องหมายดังนี้

\begin{itemize}
    \item $v_o>0$ เมื่อผู้สังเกตเคลื่อนที่เข้าหาแหล่งกำเนิด
    \item $v_o<0$ เมื่อผู้สังเกตเคลื่อนที่ออกจากแหล่งกำเนิด
    \item $v_s>0$ เมื่อแหล่งกำเนิดเคลื่อนที่เข้าหาผู้สังเกต
    \item $v_s<0$ เมื่อแหล่งกำเนิดเคลื่อนที่ออกจากผู้สังเกต
\end{itemize}
\end{tcolorbox}

\begin{center}
\begin{tikzpicture}[x=0.9cm,y=0.9cm]
    \draw[->,thick] (-0.5,0) -- (12.5,0) node[right] {$x$};
    \fill (3,0) circle (3pt) node[above=5pt]{แหล่งกำเนิด};
    \draw[->,very thick] (3,0.7) -- (4.7,0.7) node[midway,above] {$v_s$};
    \foreach \x in {4.5,6.0,7.5,9.0}{
        \draw[blue!60!black,thick] (\x,-0.8) -- (\x,0.8);
    }
    \fill (10.5,0) circle (3pt) node[above=5pt]{ผู้สังเกต};
    \draw[->,very thick] (10.5,-0.7) -- (9.0,-0.7) node[midway,below] {$v_o$};
    \draw[<->] (6.0,1.2) -- (7.5,1.2) node[midway,above] {$\lambda_{\rm front}$};
\end{tikzpicture}
\end{center}

\subsection{กรณีสำคัญที่ควรจำจากเหตุผล ไม่ใช่ท่องแยก}

\begin{enumerate}
    \item แหล่งกำเนิดอยู่นิ่ง ผู้สังเกตเคลื่อนที่
    \[
    \boxed{f_o=f_s\frac{c_s+v_o}{c_s}}
    \]

    \item ผู้สังเกตอยู่นิ่ง แหล่งกำเนิดเคลื่อนที่
    \[
    \boxed{f_o=f_s\frac{c_s}{c_s-v_s}}
    \]

    \item ทั้งสองเคลื่อนที่ ใช้สมการรวมและใส่เครื่องหมายจากทิศจริง
\end{enumerate}

\begin{tcolorbox}[title={ข้อควรระวังเมื่อมีลม}]
$c_s$ เป็นอัตราเร็วเสียงเทียบกับอากาศ ดังนั้น $v_s$ และ $v_o$ ในสมการดอปเพลอร์ต้องเป็นความเร็วเทียบกับอากาศด้วย

ถ้าโจทย์ให้ความเร็วเทียบพื้นและความเร็วลม ต้องเปลี่ยนกรอบอ้างอิงเป็นกรอบของอากาศก่อน
\end{tcolorbox}

\subsection{การสะท้อนจากวัตถุที่กำลังเคลื่อนที่}

เมื่อเสียงสะท้อนจากวัตถุเคลื่อนที่ ความถี่เปลี่ยนสองครั้ง

\begin{enumerate}
    \item วัตถุเคลื่อนที่ทำหน้าที่เป็นผู้สังเกตและรับความถี่ครั้งแรก
    \item หลังสะท้อน วัตถุนั้นทำหน้าที่เสมือนแหล่งกำเนิดที่กำลังเคลื่อนที่
\end{enumerate}

ถ้าแหล่งกำเนิดและผู้รับอยู่นิ่งร่วมกัน และตัวสะท้อนเคลื่อนที่เข้าหาด้วยอัตราเร็ว $u$

\[
f_1=f_0\frac{c_s+u}{c_s}
\]

หลังสะท้อน ตัวสะท้อนเป็นแหล่งกำเนิดเคลื่อนที่เข้าหาผู้รับ

\[
f_{\rm echo}=f_1\frac{c_s}{c_s-u}
\]

จึงได้

\begin{tcolorbox}[title={ดอปเพลอร์สองครั้งจากตัวสะท้อนที่เคลื่อนที่เข้าหา}]
\[
\boxed{
f_{\rm echo}
=f_0\frac{c_s+u}{c_s-u}
}
\]
\end{tcolorbox}

ถ้าผู้รับได้ยินทั้งเสียงตรงและเสียงสะท้อนพร้อมกัน ความถี่บีตคือ

\[
\boxed{f_{\rm beat}=|f_{\rm echo}-f_0|}
\]

% =========================================================
\section{คลื่นกระแทกและโซนิกบูม}

\subsection{จากหน้าคลื่นวงกลมสู่กรวยมัค}

ถ้าแหล่งกำเนิดเคลื่อนที่ช้ากว่าเสียง หน้าคลื่นแต่ละลูกแผ่ออกไปข้างหน้าได้ และผู้สังเกตด้านหน้าจะได้ยินเสียงก่อนแหล่งกำเนิดมาถึง

ถ้าแหล่งกำเนิดเคลื่อนที่เร็วกว่าเสียง แหล่งกำเนิดเคลื่อนที่แซงหน้าคลื่นที่ตนเองสร้าง หน้าคลื่นจำนวนมากจึงซ้อนกันบนผิวกรวย เรียกว่า \textbf{กรวยมัค} ขอบกรวยเป็นบริเวณที่ความดันเปลี่ยนอย่างรวดเร็วและผู้สังเกตได้ยินเป็น \textbf{โซนิกบูม}

นิยามเลขมัคเป็น

\[
\boxed{M\equiv\frac{v}{c_s}}
\]

โดย $M>1$ หมายถึงการเคลื่อนที่เหนือเสียง

ให้แหล่งกำเนิดเคลื่อนที่เป็นเวลา $t$ ได้ระยะ $vt$ ขณะที่คลื่นที่ปล่อยตอนเริ่มต้นมีรัศมี $c_st$ จากเรขาคณิตของเส้นสัมผัสกรวย

\[
\sin\mu=\frac{c_st}{vt}
\]

จึงได้

\begin{tcolorbox}[title={มุมมัค}]
\[
\boxed{
\sin\mu=\frac{c_s}{v}=\frac{1}{M}
}
\]
\end{tcolorbox}

\begin{center}
\begin{tikzpicture}[scale=0.78]
    \coordinate (O) at (0,0);
    \coordinate (S) at (8,0);
    \coordinate (Tup) at (2,3.464);
    \coordinate (Tdn) at (2,-3.464);
    \draw[->,thick] (-4.6,0) -- (9.6,0) node[right]{ทิศการเคลื่อนที่};
    \draw[blue!55!black,thick] (O) circle (4.0);
    \draw[red!65!black,very thick] (S) -- (Tup);
    \draw[red!65!black,very thick] (S) -- (Tdn);
    \fill (S) circle (3pt) node[above=5pt]{แหล่งกำเนิดปัจจุบัน};
    \fill (O) circle (2pt) node[below left]{จุดเริ่มต้น};
    \draw[->,very thick] (7.0,1.0) -- (8.8,1.0) node[midway,above] {$v$};
    \draw (6.8,0) arc[start angle=180,end angle=150,radius=1.2];
    \node at (6.75,0.48) {$\mu$};
    \draw[<->] (0,-4.55) -- (8,-4.55) node[midway,below] {$vt$};
    \draw[<->] (-0.35,0) -- (-0.35,4) node[midway,left] {$c_st$};
    \node[blue!55!black] at (-2.8,3.35) {หน้าคลื่นแรก};
    \node[red!65!black,rotate=-30] at (4.6,2.25) {แนวกรวยมัค};
\end{tikzpicture}
\end{center}

\subsection{เวลาที่ผู้สังเกตได้ยินโซนิกบูม}

เครื่องบินบินในแนวระดับด้วยอัตราเร็ว $v>c_s$ ที่ความสูง $h$ เมื่อกรวยมัคตัดผู้สังเกตบนพื้น เครื่องบินอยู่เลยผู้สังเกตไปข้างหน้าเป็นระยะ $x$

จากเรขาคณิต

\[
\tan\mu=\frac{h}{x}
\]

ดังนั้น

\[
x=h\cot\mu
\]

เวลาหลังจากเครื่องบินผ่านเหนือศีรษะจนได้ยินบูมคือ

\[
\Delta t=\frac{x}{v}
\]

จึงได้

\begin{tcolorbox}[title={ระยะและเวลาของโซนิกบูมสำหรับแนวบินระดับ}]
\[
\boxed{x=h\cot\mu}
\]

\[
\boxed{
\Delta t
=\frac{h}{v\tan\mu}
=\frac{h\sqrt{v^2-c_s^2}}{vc_s}
}
\]
\end{tcolorbox}

สูตรนี้เป็นผลเชิงเรขาคณิตของกรวยมัค ไม่ต้องใช้สมการความดันของคลื่นกระแทก

% =========================================================
\section{บีตและการสั่นพ้องที่ใช้ร่วมกับโจทย์เสียง}

\subsection{บีต}

เมื่อเสียงสองความถี่ใกล้กัน $f_1$ และ $f_2$ ซ้อนกัน ความดังจะเพิ่มและลดเป็นจังหวะด้วยความถี่

\[
\boxed{f_{\rm beat}=|f_1-f_2|}
\]

โจทย์ดอปเพลอร์มักให้ผู้สังเกตได้ยินเสียงตรงกับเสียงสะท้อนพร้อมกัน แล้วให้ความถี่บีตเพื่อย้อนหาอัตราเร็วของวัตถุ

\subsection{ท่อปลายปิดและค่าปรับปลายท่อ}

สำหรับท่อปิดหนึ่งปลาย ความยาวเชิงผลคือ

\[
L_{\rm eff}=L+e
\]

โดย $e$ เป็นค่าปรับปลายท่อ และเงื่อนไขการสั่นพ้องคือ

\[
\boxed{
L+e=\frac{(2n-1)\lambda}{4}
\qquad n=1,2,3,\ldots
}
\]

ตำแหน่งสั่นพ้องที่ติดกันจึงต่างกัน

\[
\boxed{L_{n+1}-L_n=\frac{\lambda}{2}}
\]

ผลนี้ช่วยให้หา $\lambda$ ได้โดยไม่ต้องทราบ $e$ ก่อน

% =========================================================
\section{ตัวอย่างประยุกต์ก่อนทำข้อสอบ}

\begin{examplebox}{ตัวอย่างที่ 1: รถพยาบาลวิ่งผ่านผู้สังเกต}
รถพยาบาลเปิดไซเรนความถี่ $700\,\si{Hz}$ เคลื่อนที่ด้วยอัตราเร็ว $25\,\si{m/s}$ อัตราเร็วเสียงเท่ากับ $350\,\si{m/s}$ ผู้สังเกตอยู่นิ่ง จงหาความถี่ก่อนและหลังรถวิ่งผ่าน

ก่อนวิ่งผ่าน แหล่งกำเนิดเคลื่อนที่เข้าหา

\[
f_{\rm approach}
=700\frac{350}{350-25}
\]

\[
\boxed{f_{\rm approach}\approx754\,\si{Hz}}
\]

หลังวิ่งผ่าน แหล่งกำเนิดเคลื่อนที่ออกจากผู้สังเกต

\[
f_{\rm recede}
=700\frac{350}{350+25}
\]

\[
\boxed{f_{\rm recede}\approx653\,\si{Hz}}
\]
\end{examplebox}

\begin{examplebox}{ตัวอย่างที่ 2: เสียงสะท้อนจากรถที่เคลื่อนที่}
เครื่องส่งเสียงอยู่นิ่งปล่อยความถี่ $800\,\si{Hz}$ รถเคลื่อนที่เข้าหาเครื่องส่งด้วยอัตราเร็ว $20\,\si{m/s}$ ให้ $c_s=340\,\si{m/s}$ จงหาความถี่ของเสียงสะท้อนที่กลับมาถึงเครื่องส่ง

\[
f_{\rm echo}
=f_0\frac{c_s+u}{c_s-u}
\]

\[
f_{\rm echo}
=800\frac{340+20}{340-20}
\]

\[
\boxed{f_{\rm echo}=900\,\si{Hz}}
\]

ถ้าได้ยินเสียงตรงพร้อมเสียงสะท้อน จะเกิดบีต

\[
\boxed{f_{\rm beat}=900-800=100\,\si{Hz}}
\]
\end{examplebox}

\begin{examplebox}{ตัวอย่างที่ 3: โซนิกบูมจากเครื่องบินมัค 2}
เครื่องบินบินด้วย $v=680\,\si{m/s}$ ที่ความสูง $h=3.00\,\si{km}$ ให้ $c_s=340\,\si{m/s}$ จงหามุมมัคและเวลาหลังผ่านเหนือศีรษะที่ผู้สังเกตได้ยินบูม

\[
M=\frac{680}{340}=2
\]

\[
\sin\mu=\frac{1}{2}
\]

\[
\boxed{\mu=30^\circ}
\]

ระยะที่เครื่องบินอยู่เลยผู้สังเกต

\[
x=h\cot\mu
=(3000)\cot30^\circ
\]

\[
x\approx5.20\times10^3\,\si{m}
\]

ดังนั้น

\[
\Delta t=\frac{x}{v}
\]

\[
\boxed{\Delta t\approx7.64\,\si{s}}
\]
\end{examplebox}

% =========================================================
\section{สรุปสูตรสำคัญ}

\begin{tcolorbox}[title={คลื่นและความเข้มเสียง}]
\[
\boxed{T=\frac{1}{f}}
\qquad
\boxed{c_s=f\lambda}
\]

\[
\boxed{I=\frac{\mathcal P}{A}}
\qquad
\boxed{I=\frac{\mathcal P}{4\pi r^2}}
\]

\[
\boxed{\frac{I_1}{I_2}=\left(\frac{r_2}{r_1}\right)^2}
\]
\end{tcolorbox}

\begin{tcolorbox}[title={ดอปเพลอร์ บีต และเสียงสะท้อนจากวัตถุเคลื่อนที่}]
\[
\boxed{f_o=f_s\frac{c_s+v_o}{c_s-v_s}}
\]

\[
\boxed{f_{\rm beat}=|f_1-f_2|}
\]

\[
\boxed{
f_{\rm echo}=f_0\frac{c_s+u}{c_s-u}
}
\]
\end{tcolorbox}

\begin{tcolorbox}[title={โซนิกบูมและกรวยมัค}]
\[
\boxed{M=\frac{v}{c_s}}
\]

\[
\boxed{\sin\mu=\frac{1}{M}=\frac{c_s}{v}}
\]

\[
\boxed{x=h\cot\mu}
\]

\[
\boxed{
\Delta t
=\frac{h}{v\tan\mu}
=\frac{h\sqrt{v^2-c_s^2}}{vc_s}
}
\]
\end{tcolorbox}

\begin{tcolorbox}[title={ท่อปิดหนึ่งปลาย}]
\[
\boxed{L+e=\frac{(2n-1)\lambda}{4}}
\]

\[
\boxed{L_{n+1}-L_n=\frac{\lambda}{2}}
\]
\end{tcolorbox}

% =========================================================
\section{ชุดข้อสอบฝึกคัดเลือกเข้า สอวน.}

\begin{tcolorbox}[title={คำชี้แจง}]
ชุดนี้เป็นข้อสอบฝึกที่เรียบเรียงใหม่ ไม่ใช่ข้อสอบ สอวน. จริง โจทย์เน้นแนวคิดระดับข้อสอบคัดเลือกเข้าค่าย 1 ได้แก่ ความสัมพันธ์ของคลื่น ความเข้มเสียง ดอปเพลอร์ บีต การสะท้อนจากวัตถุเคลื่อนที่ การสั่นพ้อง และเรขาคณิตของโซนิกบูม

โจทย์ตั้งใจให้แนวคิดพื้นฐานไม่เกินหลักสูตร แต่แต่ละข้อต้องเชื่อมอย่างน้อยสองขั้น คะแนนรวม $100$ คะแนน และมีโจทย์ท้าทายเพิ่มเติม $10$ คะแนน
\end{tcolorbox}

% ---------------------------------------------------------
\begin{exercisebox}{ข้อ 1: ผู้สังเกตเคลื่อนที่ระหว่างแหล่งกำเนิดสองแหล่ง $10$ คะแนน}
ลำโพง $A$ และ $B$ อยู่นิ่งและอยู่คนละด้านของนักเรียน ทั้งสองปล่อยเสียงความถี่เท่ากัน $510\,\si{Hz}$ นักเรียนวิ่งจาก $A$ ไปหา $B$ ด้วยอัตราเร็ว $5.0\,\si{m/s}$ ให้ $c_s=340\,\si{m/s}$

\begin{enumerate}[label=(\alph*)]
    \item หาความถี่ที่ได้ยินจากลำโพงแต่ละตัว
    \item หาความถี่บีตที่นักเรียนได้ยิน
\end{enumerate}
\end{exercisebox}

% ---------------------------------------------------------
\begin{exercisebox}{ข้อ 2: หาอัตราเร็วและความถี่จริงของไซเรน $10$ คะแนน}
รถคันหนึ่งเปิดไซเรนและวิ่งเป็นเส้นตรงผ่านผู้สังเกตที่อยู่นิ่ง ก่อนรถผ่าน ผู้สังเกตได้ยิน $660\,\si{Hz}$ และหลังรถผ่านได้ยิน $550\,\si{Hz}$ ให้ $c_s=330\,\si{m/s}$

จงหา

\begin{enumerate}[label=(\alph*)]
    \item อัตราเร็วของรถ
    \item ความถี่จริงของไซเรน
\end{enumerate}
\end{exercisebox}

% ---------------------------------------------------------
\begin{exercisebox}{ข้อ 3: ดอปเพลอร์เมื่อมีลม $12$ คะแนน}
ลมพัดไปทางทิศตะวันออกด้วยอัตราเร็ว $15\,\si{m/s}$ รถที่เป็นแหล่งกำเนิดอยู่ด้านตะวันตกและเคลื่อนที่ไปทางตะวันออกด้วยอัตราเร็วเทียบพื้น $45\,\si{m/s}$ ผู้สังเกตอยู่ด้านตะวันออกและเคลื่อนที่ไปทางตะวันออกด้วยอัตราเร็วเทียบพื้น $5\,\si{m/s}$ แหล่งกำเนิดปล่อยเสียง $600\,\si{Hz}$ อัตราเร็วเสียงเทียบกับอากาศคือ $340\,\si{m/s}$

\begin{enumerate}[label=(\alph*)]
    \item หาความเร็วของแหล่งกำเนิดและผู้สังเกตเทียบกับอากาศ พร้อมระบุว่ากำลังเข้าหาหรือออกจากกันในกรอบอากาศ
    \item หาความถี่ที่ผู้สังเกตได้รับ
\end{enumerate}
\end{exercisebox}

% ---------------------------------------------------------
\begin{exercisebox}{ข้อ 4: ไซเรนหมุนเป็นวงกลม $10$ คะแนน}
ไซเรนติดอยู่ที่ปลายแขนหมุนรัศมี $1.50\,\si{m}$ ผู้สังเกตอยู่ไกลมากเมื่อเทียบกับรัศมีวงกลม จึงประมาณได้ว่าองค์ประกอบความเร็วของไซเรนตามแนวสายตาสูงสุดมีขนาดเท่ากับอัตราเร็วตามเส้นรอบวง

ผู้สังเกตวัดความถี่สูงสุดและต่ำสุดได้ $660\,\si{Hz}$ และ $440\,\si{Hz}$ ตามลำดับ ให้ $c_s=330\,\si{m/s}$

\begin{enumerate}[label=(\alph*)]
    \item หาอัตราเร็วตามเส้นรอบวงของไซเรน
    \item หาความถี่จริงของไซเรน
    \item หาอัตราเร็วเชิงมุมของแขนหมุน
\end{enumerate}
\end{exercisebox}

% ---------------------------------------------------------
\begin{exercisebox}{ข้อ 5: บีตจากรถที่สะท้อนเสียง $12$ คะแนน}
เครื่องส่งและเครื่องรับเสียงอยู่ตำแหน่งเดียวกันและอยู่นิ่ง เครื่องส่งปล่อยเสียงความถี่ $1000\,\si{Hz}$ ไปยังรถที่กำลังเคลื่อนที่เข้าหา เสียงสะท้อนกลับมาซ้อนกับเสียงตรง ทำให้ได้ยินบีต $125\,\si{Hz}$ ให้ $c_s=340\,\si{m/s}$

\begin{enumerate}[label=(\alph*)]
    \item อธิบายว่าความถี่เปลี่ยนสองครั้งอย่างไร
    \item หาอัตราเร็วของรถ
\end{enumerate}
\end{exercisebox}

% ---------------------------------------------------------
\begin{exercisebox}{ข้อ 6: หาค่าปรับปลายท่อจากการสั่นพ้อง $10$ คะแนน}
ส้อมเสียงความถี่ $500\,\si{Hz}$ วางเหนือท่อเรโซแนนซ์ที่ปิดด้านล่างด้วยผิวน้ำ พบการสั่นพ้องสองตำแหน่งติดกันเมื่อความยาวลำอากาศเป็น $0.150\,\si{m}$ และ $0.490\,\si{m}$

\begin{enumerate}[label=(\alph*)]
    \item หาความยาวคลื่นและอัตราเร็วเสียง
    \item หาค่าปรับปลายท่อ $e$
    \item ถ้าใช้แบบจำลอง $e=0.60r$ จงหารัศมีท่อ $r$
\end{enumerate}
\end{exercisebox}

% ---------------------------------------------------------
\begin{exercisebox}{ข้อ 7: รถตำรวจแซงรถยนต์ $10$ คะแนน}
รถตำรวจเปิดไซเรนความถี่คงตัวและวิ่งไปในทิศเดียวกับรถยนต์คันหนึ่ง รถยนต์มีอัตราเร็ว $10.0\,\si{m/s}$ ขณะที่รถตำรวจอยู่ด้านหลังและกำลังไล่เข้าใกล้ คนขับรถยนต์ได้ยินความถี่ $720\,\si{Hz}$ หลังจากรถตำรวจแซงไปแล้ว คนขับได้ยินความถี่ $600\,\si{Hz}$

ให้ความเร็วเสียง $c_s=340\,\si{m/s}$ และไม่มีลม จงหา

\begin{enumerate}[label=(\alph*)]
    \item อัตราเร็วของรถตำรวจ
    \item ความถี่จริงของไซเรน
\end{enumerate}
\end{exercisebox}

% ---------------------------------------------------------
\begin{exercisebox}{ข้อ 8: ไซเรนตกอิสระ $8$ คะแนน}
ไซเรนความถี่ $600\,\si{Hz}$ ถูกปล่อยตกจากหยุดนิ่งในแนวดิ่งเหนือผู้สังเกต ก่อนถึงผู้สังเกต ไซเรนตกมาแล้วเป็นระยะ $45\,\si{m}$ ให้ $g=10\,\si{m/s^2}$ และ $c_s=330\,\si{m/s}$

โดยละเลยแรงต้านอากาศ จงหาความถี่ที่ผู้สังเกตได้รับจากเสียงที่ปล่อย ณ ขณะนั้น
\end{exercisebox}

% ---------------------------------------------------------
\begin{exercisebox}{ข้อ 9: เครื่องบินมัค 2 $10$ คะแนน}
เครื่องบินบินในแนวระดับด้วยอัตราเร็ว $680\,\si{m/s}$ ที่ความสูง $3.00\,\si{km}$ ให้ $c_s=340\,\si{m/s}$

\begin{enumerate}[label=(\alph*)]
    \item หาเลขมัคและมุมมัค
    \item เมื่อผู้สังเกตได้ยินโซนิกบูม เครื่องบินอยู่เลยจุดเหนือศีรษะไปตามแนวราบเป็นระยะเท่าใด
    \item ผู้สังเกตได้ยินบูมหลังเครื่องบินผ่านเหนือศีรษะกี่วินาที
\end{enumerate}
\end{exercisebox}

% ---------------------------------------------------------
\begin{exercisebox}{ข้อ 10: หาอัตราเร็วและความสูงจากผู้สังเกตสองคน $8$ คะแนน}
ผู้สังเกต $A$ และ $B$ อยู่บนแนวเส้นทางบินเดียวกัน โดย $B$ อยู่ข้างหน้า $A$ เป็นระยะ $3.00\,\si{km}$ เครื่องบินเหนือเสียงบินในแนวระดับผ่านเหนือทั้งสองคนด้วยอัตราเร็วคงตัว ทั้งสองได้ยินโซนิกบูมห่างกัน $5.00\,\si{s}$

ผู้สังเกต $A$ เห็นเครื่องบินผ่านเหนือศีรษะ แล้วได้ยินบูมหลังจากนั้น $8.00\,\si{s}$ ให้ $c_s=340\,\si{m/s}$

\begin{enumerate}[label=(\alph*)]
    \item หาอัตราเร็วของเครื่องบิน
    \item หาเลขมัคและมุมมัค
    \item หาความสูงของเครื่องบิน
\end{enumerate}
\end{exercisebox}

% ---------------------------------------------------------
\begin{exercisebox}{โจทย์ท้าทายเพิ่มเติม: รถไฟสองขบวนและลม $10$ คะแนน}
ลมพัดไปทางตะวันออกด้วยอัตราเร็ว $10\,\si{m/s}$ รถไฟ $A$ อยู่ทางตะวันตก เคลื่อนที่ไปทางตะวันออกด้วยอัตราเร็วเทียบพื้น $20\,\si{m/s}$ และมีเครื่องส่งกับเครื่องรับเสียงอยู่บนขบวนเดียวกัน รถไฟ $B$ อยู่ทางตะวันออก เคลื่อนที่ไปทางตะวันตกด้วยอัตราเร็วเทียบพื้น $10\,\si{m/s}$ และสะท้อนเสียงได้ดี

รถไฟ $A$ ส่งเสียงความถี่ $800\,\si{Hz}$ ไปยังรถไฟ $B$ ให้ $c_s=340\,\si{m/s}$ เทียบกับอากาศ จงหาความถี่ของเสียงสะท้อนที่กลับมาถึงรถไฟ $A$
\end{exercisebox}

% =========================================================
\section{เฉลยละเอียดและการตรวจคำตอบ}

\subsection*{เฉลยข้อ 1}

ลำโพง $B$ อยู่ด้านหน้า นักเรียนเคลื่อนที่เข้าหา

\[
f_B=510\frac{340+5}{340}
\]

\[
\boxed{f_B=517.5\,\si{Hz}}
\]

ลำโพง $A$ อยู่ด้านหลัง นักเรียนเคลื่อนที่ออกจาก

\[
f_A=510\frac{340-5}{340}
\]

\[
\boxed{f_A=502.5\,\si{Hz}}
\]

ดังนั้น

\[
f_{\rm beat}=|517.5-502.5|
\]

\[
\boxed{f_{\rm beat}=15.0\,\si{Hz}}
\]

% ---------------------------------------------------------
\subsection*{เฉลยข้อ 2}

ก่อนผ่าน

\[
660=f_s\frac{330}{330-v}
\tag{1}
\]

หลังผ่าน

\[
550=f_s\frac{330}{330+v}
\tag{2}
\]

หาร (1) ด้วย (2)

\[
\frac{660}{550}=\frac{330+v}{330-v}
\]

จัดรูปทั่วไปได้

\[
v=c_s\frac{f_{\rm approach}-f_{\rm recede}}
{f_{\rm approach}+f_{\rm recede}}
\]

แทนค่า

\[
v=330\frac{660-550}{660+550}
\]

\[
\boxed{v=30.0\,\si{m/s}}
\]

แทนกลับใน (1)

\[
f_s=660\frac{330-30}{330}
\]

\[
\boxed{f_s=600\,\si{Hz}}
\]

% ---------------------------------------------------------
\subsection*{เฉลยข้อ 3}

เลือกทิศตะวันออกเป็นบวก ความเร็วเทียบอากาศเท่ากับความเร็วเทียบพื้นลบความเร็วลม

แหล่งกำเนิด

\[
v_{s,\rm air}=45-15=30\,\si{m/s}
\]

จึงเคลื่อนที่เข้าหาผู้สังเกตด้วย $30\,\si{m/s}$

ผู้สังเกต

\[
v_{o,\rm air}=5-15=-10\,\si{m/s}
\]

เครื่องหมายลบในแกนตะวันออกหมายถึงผู้สังเกตเคลื่อนที่ไปทางตะวันตกในกรอบอากาศ ซึ่งเป็นการเคลื่อนที่เข้าหาแหล่งกำเนิด จึงใช้ขนาด $v_o=10\,\si{m/s}$ ในสูตรเครื่องหมายแบบเข้าหา

\[
f_o=600\frac{340+10}{340-30}
\]

\[
\boxed{f_o\approx677.4\,\si{Hz}}
\]

ตรวจแนวโน้ม: ทั้งแหล่งกำเนิดและผู้สังเกตเคลื่อนที่เข้าหากันในกรอบอากาศ ดังนั้นความถี่ต้องสูงกว่า $600\,\si{Hz}$

% ---------------------------------------------------------
\subsection*{เฉลยข้อ 4}

ที่ความถี่สูงสุด แหล่งกำเนิดมีองค์ประกอบความเร็วเข้าหาผู้สังเกตขนาด $v$

\[
660=f_s\frac{330}{330-v}
\tag{1}
\]

ที่ความถี่ต่ำสุด แหล่งกำเนิดเคลื่อนที่ออก

\[
440=f_s\frac{330}{330+v}
\tag{2}
\]

จึงได้

\[
v=330\frac{660-440}{660+440}
\]

\[
\boxed{v=66.0\,\si{m/s}}
\]

ความถี่จริง

\[
f_s=\frac{2f_{\max}f_{\min}}{f_{\max}+f_{\min}}
\]

\[
\boxed{f_s=528\,\si{Hz}}
\]

อัตราเร็วเชิงมุม

\[
\omega=\frac{v}{R}
=\frac{66.0}{1.50}
\]

\[
\boxed{\omega=44.0\,\si{rad/s}}
\]

% ---------------------------------------------------------
\subsection*{เฉลยข้อ 5}

รถรับเสียงในฐานะผู้สังเกตที่เคลื่อนที่เข้าหา

\[
f_1=f_0\frac{c_s+u}{c_s}
\]

จากนั้นรถสะท้อนเสียงและทำหน้าที่เป็นแหล่งกำเนิดที่เคลื่อนที่เข้าหาเครื่องรับ

\[
f_{\rm echo}=f_1\frac{c_s}{c_s-u}
=f_0\frac{c_s+u}{c_s-u}
\]

รถเข้าหา จึงมี $f_{\rm echo}>f_0$ และ

\[
f_{\rm echo}=1000+125=1125\,\si{Hz}
\]

ดังนั้น

\[
\frac{1125}{1000}=\frac{340+u}{340-u}
\]

\[
1.125(340-u)=340+u
\]

\[
42.5=2.125u
\]

\[
\boxed{u=20.0\,\si{m/s}}
\]

% ---------------------------------------------------------
\subsection*{เฉลยข้อ 6}

ตำแหน่งสั่นพ้องติดกันของท่อปิดหนึ่งปลายต่างกัน $\lambda/2$

\[
L_2-L_1=\frac{\lambda}{2}
\]

\[
\lambda=2(0.490-0.150)
\]

\[
\boxed{\lambda=0.680\,\si{m}}
\]

\[
c_s=f\lambda=(500)(0.680)
\]

\[
\boxed{c_s=340\,\si{m/s}}
\]

ตำแหน่งแรกมี

\[
L_1+e=\frac{\lambda}{4}
\]

\[
e=\frac{0.680}{4}-0.150
\]

\[
\boxed{e=0.0200\,\si{m}}
\]

จาก $e=0.60r$

\[
r=\frac{0.0200}{0.60}
\]

\[
\boxed{r=3.33\times10^{-2}\,\si{m}=3.33\,\si{cm}}
\]

% ---------------------------------------------------------
\subsection*{เฉลยข้อ 7}

ให้รถตำรวจมีอัตราเร็ว $v_s$ และความถี่จริงของไซเรนเป็น $f_s$

ก่อนแซง เสียงเดินทางไปข้างหน้า รถยนต์เคลื่อนที่หนีหน้าคลื่น แต่รถตำรวจเคลื่อนที่เข้าหารถยนต์ จึงมี

\[
720
=f_s\frac{340-10}{340-v_s}
=f_s\frac{330}{340-v_s}
\tag{1}
\]

หลังแซง เสียงที่ถึงรถยนต์เดินทางย้อนกลับ รถยนต์เคลื่อนที่เข้าหาหน้าคลื่น ส่วนรถตำรวจเคลื่อนที่ออกจากรถยนต์ จึงมี

\[
600
=f_s\frac{340+10}{340+v_s}
=f_s\frac{350}{340+v_s}
\tag{2}
\]

หารสมการ (1) ด้วยสมการ (2)

\[
\frac{720}{600}
=\frac{330(340+v_s)}{350(340-v_s)}
\]

จัดรูป

\[
420(340-v_s)=330(340+v_s)
\]

\[
142800-420v_s=112200+330v_s
\]

\[
30600=750v_s
\]

\[
\boxed{v_s=40.8\,\si{m/s}}
\]

แทนในสมการ (1)

\[
f_s
=720\frac{340-40.8}{330}
\]

\[
\boxed{f_s=652.8\,\si{Hz}\approx653\,\si{Hz}}
\]

ตรวจสอบกับสมการหลังแซง

\[
652.8\frac{350}{340+40.8}=600\,\si{Hz}
\]

จึงสอดคล้องกับข้อมูลทั้งสองช่วง

% ---------------------------------------------------------
\subsection*{เฉลยข้อ 8}

จากการตกอิสระ

\[
v^2=2gh
\]

\[
v=\sqrt{2(10)(45)}
\]

\[
\boxed{v=30.0\,\si{m/s}}
\]

ไซเรนเคลื่อนที่เข้าหาผู้สังเกตที่อยู่นิ่ง

\[
f_o=600\frac{330}{330-30}
\]

\[
\boxed{f_o=660\,\si{Hz}}
\]

% ---------------------------------------------------------
\subsection*{เฉลยข้อ 9}

\[
M=\frac{680}{340}=2.00
\]

\[
\sin\mu=\frac{1}{M}=\frac{1}{2}
\]

\[
\boxed{\mu=30.0^\circ}
\]

ระยะตามแนวราบ

\[
x=h\cot\mu
=(3000)\cot30^\circ
\]

\[
\boxed{x\approx5.20\times10^3\,\si{m}}
\]

เวลา

\[
\Delta t=\frac{x}{v}
=\frac{5.196\times10^3}{680}
\]

\[
\boxed{\Delta t\approx7.64\,\si{s}}
\]

% ---------------------------------------------------------
\subsection*{เฉลยข้อ 10}

สำหรับผู้สังเกตทั้งสองที่อยู่ใต้แนวบินและเครื่องบินมีความสูงคงตัว เวลาหน่วงเชิงเรขาคณิตหลังเครื่องบินผ่านเหนือผู้สังเกตแต่ละคนเท่ากัน ดังนั้นผลต่างเวลาที่ได้ยินบูมเท่ากับเวลาที่เครื่องบินเดินทางจาก $A$ ไป $B$

\[
v=\frac{3000}{5.00}
\]

\[
\boxed{v=600\,\si{m/s}}
\]

\[
M=\frac{600}{340}
\]

\[
\boxed{M\approx1.765}
\]

\[
\sin\mu=\frac{340}{600}
\]

\[
\boxed{\mu\approx34.52^\circ}
\]

จาก

\[
\Delta t=\frac{h}{v\tan\mu}
\]

\[
h=\Delta t\,v\tan\mu
\]

\[
h=(8.00)(600)\tan34.52^\circ
\]

\[
\boxed{h\approx3.30\times10^3\,\si{m}}
\]

% ---------------------------------------------------------
\subsection*{เฉลยโจทย์ท้าทาย}

เลือกกรอบอากาศ ความเร็วเทียบอากาศของรถไฟ $A$ คือ

\[
v_A=20-10=10\,\si{m/s}
\]

ไปทางตะวันออก ส่วนรถไฟ $B$

\[
v_B=-10-10=-20\,\si{m/s}
\]

คือไปทางตะวันตกด้วย $20\,\si{m/s}$ ในกรอบอากาศ

ช่วงขาไป รถไฟ $A$ เป็นแหล่งกำเนิดเคลื่อนที่เข้าหา $B$ ด้วย $10\,\si{m/s}$ และ $B$ เป็นผู้สังเกตเคลื่อนที่เข้าหาเสียงด้วย $20\,\si{m/s}$

\[
f_B
=800\frac{340+20}{340-10}
\]

\[
f_B\approx872.73\,\si{Hz}
\]

ช่วงขากลับ รถไฟ $B$ ทำหน้าที่เป็นแหล่งกำเนิดเคลื่อนที่เข้าหา $A$ ด้วย $20\,\si{m/s}$ และรถไฟ $A$ เคลื่อนที่เข้าหาเสียงขากลับด้วย $10\,\si{m/s}$

\[
f_A
=f_B\frac{340+10}{340-20}
\]

\[
f_A
=800
\frac{340+20}{340-10}
\frac{340+10}{340-20}
\]

\[
\boxed{f_A\approx954.5\,\si{Hz}}
\]



\fussy